\pg{424} \emph{Anthropologia}, a visão \emph{intuitiva} de \textsc{Deus} excede a natural capacidade da nossa intelligencia. Logo á existencia de \textsc{Deus} deve ser demonstrada \footnote{Os que negam a necessidade de demonstrar a existencia de \textsc{Deus} são --- a) os \emph{sentimentalistas}, sustentando que conhecemos \emph{immediatamente} a \textsc{DEus} por meio de um sentido arcano e divino. --- 5) os \emph{ontologistas}, admittindo uma \emph{zntuicüo immediata} ou \emph{directa} de \textsc{Deus}.\par Estas duas opiniões são falsas. -\par \emph{a}) E' falsa a dos \emph{sentimentalistas}; pois a existencia d'esse sentido arcano e divino é contraria ao testemunho da consciencia.\par 5) E' falsa a dos \emph{ontologistas}; pois o \emph{ontologismo}, como vimos, € falso.\par Aqui apenas accrescentaremos que a S. Egreja, Mestra infallivel de verdade, declarou suspeitas e perigosas, entre outras, estas proposições: « Immediata Dk: cognitio, habitualis saltem, intellectui humano essentialis est, ita ut sine ea nihil cognoscere possit: siquidem est ipsum lumen intel lectuale ». --- « Esse illud, quod in omnibus et sine quo nihil intelligimus, est esse divinum ». (Decr. S. Officii, die 18 Septemb. a. 1861).}.

8. \emph{A} \textbf{existencia} de \textbf{Deus} não se demonstra \textbf{a-priori}. --- A demonstração diz-se \emph{a-priori}, quando a existencia de um ser se prova 